\chapter{Permasalahan AI}
Pada bab ini akan kita uraikan beberapa permasalahan atau tantangan AI.
Ada masalah yang muncul karena teknologi yang tersedia saat ini masih
terbatas atau karena pendekatan (madzhab) yang digunakan kurang cocok
dengan permasalahan yang ingin dipecahkan.

\section{Halusinasi}
Sistem AI yang menggunakan mekanisme \emph{learning} seperti aliran
\emph{connectionist} membutuhkan contoh data yang sangat banyak.
Jika dataset yang digunakan kurang atau tidak ada, maka dapat terjadi 
halusinasi.
Mari kita ambil contoh.

Katakan kita sedang membuat sebuah model AI yang ingin membedakan
kucing dan anjing dari gambar (image).
Maka kita kumpulkan gambar kucing dan anjing.
Secara umum (rule of thumb) dibutuhkan 10.000 contoh gambar untuk
masing-masing kelas. Jadi kita harus mengumpulkan 10.000 gambar kucing
dan 10.000 gambar anjing. Kadang dengan 5000 gambar juga sudah dapat
membedakan tetapi akurasinya masih rendah. Apalagi jika datanya kurang
dari 5000! Hasilnya kurang bagus\footnote{Pernah kejadian ada mahasiswa
yang akurasi dari modelnya hanya 50\%. Lah ini sih sama dengan melempar
koin.}.

Katakanlah kita sudah berhasil mengumpulkan gambar sesuai dengan kebutuhan
tersebut. Sepuluh ribu data untuk masing-masing kelas.
Kemudian kita \emph{training} arsitektur (rangkaian) AI kita sehingga
menghasilkan model yang memiliki akurasi cukup baik.
Proses training inilah yang sangat membutuhkan GPU dan super lama
sehingga sangat mahal.

Model ini kemudian kita uji dengan data pengujian (test).
Katakan hasilnya sangat baik.
Bagaimana jika kita uji dengan \textbf{gambar kuda}?
Kira-kira apa hasilnya? Kucing atau anjing?
Mesin AI itu \textbf{HARUS} menjawab.
Karena dia hanya mengerti kucing atau anjing maka pilihannya hanya itu.
Kemungkinan dia akan menjawab dengan anjing karena ukurannya yang besar.
Maka dia mengalami halusinasi.

Katakan sekarang kita tambahkan 10.000 gambar kuda dan kita latih kembali
mesin AI kita, maka sekarang mesin AI kita dapat mengenali kucing, anjing,
dan kuda.
Mari kita tanya (tes) mesin AI kita dengan gambar gajah.
Dia akan kebingungan dan mungkin akan menjawab bahwa itu gambar kuda.
Halusinasi kembali.
Kita harus menambahkan gambar gajah ke dalam data latih dan melatih kembali.
Prosesnya berulang.
Dia akan kembali halusinasi jika kita tanya dengan data yang dia belum
pernah kenal.

Halusinasi terjadi karena mesin AI tersebut tidak dilatih dengan menggunakan
data yang (akan) ditanyakan. Jadi dia akan menjawab dengan apa yang
dia tahu. Dengan super percaya diri.
Itulah sebabnya kita jangan seratus persen percaya kepada jawaban AI.
Agak mengerikan juga banyak orang yang sangat percaya kepada AI.

\section{Tidak Dapat Dijelaskan}
Umumnya kita tidak tahu kenapa keluaran dari mesin AI seperti itu.
AI ini seperti kotak hitam (\emph{black box}).
Dia diberi masukan dan mengeluarkan sesuatu.
Mengapa keluarannya seperti itu, tidak jelas.

Untuk mengurangi ketidakjelasan tersebut maka ada pendekatan
\emph{explainable AI} atau xAI.
Ini akan dijelaskan pada update (versi) berikutnya.

\section{Tidak Transparan}
Mirip seperti sebelumnya, AI juga tidak transparan.
Jika kita mendapatkan sebuah model AI yang sudah jadi,
seringkali kita tidak tahu arsitektur yang digunakan,
data set yang digunakan pada proses pelatihan,
darimana data set tersebut diperoleh, dan seterusnya.

\section{Bias - Tidak Fair}
Model AI juga dapat menghasilkan keluaran atau keputusan yang tidak fair.
Ada contoh kasus di Amazon dalam hal pengelolaan sumber daya manusia (SDM).
Hal ini dapat terjadi karena data set tidak seimbang.
Ada kelas yang lebih banyak datanya dibandingkan kelas lainnya.
Sebagai contoh, untuk data SDM biasanya data laki-laki lebih banyak
daripada data perempuan. Data yang njomplang ini akan menghasilkan bias.

\section{Terlalu Banyak Menggunakan Energi}
Seperti telah dijelaskan sebelumnya, AI membutuhkan GPU dalam jumlah yang besar
(banyak).
GPU-GPU ini memang membutuhkan listrik (energi) yang banyak juga.
Ke depannya harus didesain arsitektur komputasi yang lebih hemat energi
tetapi masih memberikan kinerja (performance) yang bagus juga.

\section{AI Menghabiskan Air?}
Beredar isyu bahwa AI menghabiskan air bersih. Benarkah?
Jadi ceritanya begini.

Sebagaimana telah diuraikan sebelumnya, AI - setidaknya yang menggunakan
madzhab \emph{connectionist} - membutuhkan GPU dalam prosesnya.
GPU-GPU ini menghasilkan panas yang luar biasa. Seringkali panas chip
GPU ini melebihi panasnya CPU. Mungkin karena mereka sedang bekerja keras
terus menerus.
Akibatnya ruangan (data center) yang digunakan untuk menyimpan GPU-GPU ini
menjadi sangat panas. 
Pendinginan di data center menggunakan AC dan sekarang menggunakan air
seperti radiator di mobil. Nah ``radiator'' ini membutuhkan air bersih.
Jadi inilah yang sering disebut AI menghabiskan air.
Itulah sebabnya banyak data center yang dipergunakan untuk AI dekat dengan
sungai (sumber air).

Sesungguhnya pengelolaan pendinginan inilah yang menjadi kuncinya.
Jadi tidak benar AI itu seperti minum air.
